\documentclass[reqno]{amsart} 
\usepackage{amssymb,latexsym,amsmath,amscd,graphicx,setspace,amsthm,verbatim,enumitem}
\usepackage[margin = 3 cm]{geometry}

\def\upint{\mathchoice%
    {\mkern13mu\overline{\vphantom{\intop}\mkern7mu}\mkern-20mu}%
    {\mkern7mu\overline{\vphantom{\intop}\mkern7mu}\mkern-14mu}%
    {\mkern7mu\overline{\vphantom{\intop}\mkern7mu}\mkern-14mu}%
    {\mkern7mu\overline{\vphantom{\intop}\mkern7mu}\mkern-14mu}%
  \int}
\def\lowint{\mkern3mu\underline{\vphantom{\intop}\mkern7mu}\mkern-10mu\int}


\theoremstyle{plain}
\newtheorem{theorem}{Theorem}[section]
\newtheorem{proposition}{Proposition}
\newtheorem{corollary}{Corollary}
\newtheorem{lemma}{Lemma}
\newtheorem{conjecture}{Conjecture}
\newtheorem{question}{Question}
\newtheorem{problem}{Problem}
      
\theoremstyle{definition}
\newtheorem{definition}{Definition}

\newenvironment{solution1}{\paragraph{\emph{Solution $1$}.}}{\hfill$\square$}
\newenvironment{solution2}{\paragraph{\emph{Solution $2$}.}}{\hfill$\square$}
\newenvironment{solution3}{\paragraph{\emph{Solution $3$}.}}{\hfill$\square$}

\begin{document} 

\title[Homework 4]{Homework 4}

\date{\today} 
\maketitle 



\begin{problem}
Show that a subset $A$ of $\mathbb{R}^{n}$ is Lebesgue measurable if and only if for all $\varepsilon > 0$, there exists a closed set $F \subseteq \mathbb{R}^{n}$ such that
\begin{enumerate}
\item $F \subseteq A$,
\item $\lambda^{*}(A \smallsetminus F) < \varepsilon$.
\end{enumerate}
(Hint:  We showed in class that if $A$ is Lebesgue measurable, then so is $A^{c}$.)
\end{problem}
\begin{solution1}

\end{solution1}



\begin{problem}
Show that if $I_{1},\ldots,I_{m}$ are finitely many non-overlapping $n$-dimensional closed boxes such that
$$\bigcup_{i=1}^{m} I_{i} \subseteq \bigcup_{j=1}^{s}J_{j}, $$
where $J_{1},\ldots,J_{s}$ are finitely many $n$-dimensional closed boxes, then 
$$\sum_{i=1}^{m}V_{n}(I_{i}) \le \sum_{j=1}^{s}V_{n}(J_{j}). $$
(You can try for $n=1$ first and then prove it for $n=2$.  That would be enough.  We used that result in class.)
\end{problem}
\begin{solution1}

\end{solution1}


\begin{problem}
Show that a finite union of compact sets is a compact set.  (Try to give a proof using the original definition of compactness in terms of open covers.  We also used that result in class.)
\end{problem}
\begin{solution1}

\end{solution1}



\begin{problem}
Let $K \subseteq \mathbb{R}^{n}$ be a compact set and let $(x_{m})_{m=1}^{\infty}$ be a sequence of elements of $K$.  Show that the sequence $(x_{m})_{m=1}^{\infty}$ has a convergent subsequence that converges to an element of $K$.  (We will use this result in the upcoming week.)
\end{problem}
\begin{solution1}

\end{solution1}




Abstracting the properties that Lebesgue measurable subsets of $\mathbb{R}^{n}$ satisfy, one is led to the notion of a $\sigma$-algebra.  If $S$ is an arbitrary set, a subset $\mathcal{A} \subseteq \mathcal{P}(S)$ is called a $\sigma$-algebra if:
\begin{enumerate}
\item $\varnothing \in \mathcal{A}$,
\item Given a countable collection $A_{1},A_{2},\ldots$ of sets in $\mathcal{A}$, one has $\bigcup_{k=1}^{\infty}A_{k} \in \mathcal{A}$,
\item If $A \in \mathcal{A}$, then $A^{c} \in \mathcal{A}$.
\end{enumerate}
 
 
\begin{problem}
Let $\mathcal{A} \subseteq \mathcal{P}(S)$ be a $\sigma$-algebra.  Show that if $A_{1},A_{2},\ldots$ is a countable collection of sets in $\mathcal{A}$, then $\bigcap_{k=1}^{\infty}A_{k} \in \mathcal{A}$.
\end{problem}
\begin{solution1}

\end{solution1} 


\begin{problem}[2B]
Show that $S= \{\bigcup_{n \in K} (n, n+1]: K \subseteq \mathbb{Z} \}$ is a $\sigma$-algebra on $\mathbb{R}$.
\end{problem}
\begin{solution1}

\end{solution1}







\end{document}









